\chapter{Integration by Parts}

Integration by parts is the product rule
run backwards. It is the standard weapon against integrands that mix
unlike species: a polynomial times an exponential, a polynomial times a
logarithm, a logarithm on its own, and so on.

\section{Core Rules}

\begin{rulebox}{The Formula}
For differentiable $u$ and $v$,
\[
  \int u(x)\,v'(x)\,\mathrm{d}x
    \;=\; u(x)v(x) \;-\; \int u'(x)\,v(x)\,\mathrm{d}x .
\]
Pick one factor of the integrand to be $u$ (which you will differentiate)
and the rest to be $v'$ (which you will antidifferentiate to $v$). The
goal: $\int u'v\,\mathrm{d}x$ should be \emph{simpler} than
$\int uv'\,\mathrm{d}x$.
\end{rulebox}

\begin{rulebox}{Definite Version}
\[
  \int_{a}^{b} u(x)\,v'(x)\,\mathrm{d}x
    \;=\; \bigl[u(x)v(x)\bigr]_{a}^{b} \;-\; \int_{a}^{b} u'(x)\,v(x)\,\mathrm{d}x .
\]
\end{rulebox}

\begin{rulebox}{Choosing $u$: the LIATE Heuristic}
Prefer $u$ in the following order, pick the first type that appears:
\begin{itemize}[leftmargin=1.6em,itemsep=0.3ex]
  \item Logarithm: $\ln x$, $\log x$;
  \item Inverse trig (rarely needed here);
  \item Algebraic: polynomials, powers;
  \item Trigonometric: $\sin x$, $\cos x$;
  \item Exponential: $\mathrm{e}^{x}$, $\mathrm{e}^{ax}$.
\end{itemize}
For $\int x^{n}\mathrm{e}^{ax}\,\mathrm{d}x$ pick $u=x^{n}$ (A beats E);
for $\int x^{n}\ln x\,\mathrm{d}x$ pick $u=\ln x$ (L beats A).
\end{rulebox}

\begin{rulebox}{Repeated Application and the Tabular Method}
If a single application of parts still leaves an integrable but
compound expression, apply parts again. For
$\int x^{n}\mathrm{e}^{ax}\,\mathrm{d}x$ with integer $n\geq 1$, the
tabular shortcut is to list successive derivatives of $u=x^{n}$ in one
column (until you hit $0$) and successive antiderivatives of
$\mathrm{e}^{ax}$ in the other, then read off the answer as an
alternating sum of products along the diagonals.
\end{rulebox}

\begin{rulebox}{The Trick $\int \ln x\,\mathrm{d}x$}
Treat the integrand as $u\cdot v'$ with $u=\ln x$, $v'=1$. Then $v=x$
and parts gives
\[
  \int \ln x\,\mathrm{d}x
    = x\ln x - \int x\cdot\frac1x\,\mathrm{d}x = x\ln x - x + C .
\]
\end{rulebox}

\section{Practice Problems}

Compute each integral using parts (possibly more than once).

\begin{problems}
  \item $\displaystyle\int x\,\mathrm{e}^{x}\,\mathrm{d}x$
  \item $\displaystyle\int x\,\mathrm{e}^{-x}\,\mathrm{d}x$
  \item $\displaystyle\int x\ln x\,\mathrm{d}x$
  \item $\displaystyle\int \ln x\,\mathrm{d}x$
  \item $\displaystyle\int x^{2}\mathrm{e}^{x}\,\mathrm{d}x$
  \item $\displaystyle\int x^{2}\ln x\,\mathrm{d}x$
  \item $\displaystyle\int (\ln x)^{2}\,\mathrm{d}x$
  \item $\displaystyle\int x\,\mathrm{e}^{2x}\,\mathrm{d}x$
  \item $\displaystyle\int x(x+1)^{8}\,\mathrm{d}x$
        \quad \emph{(parts; substitution also works)}
  \item $\displaystyle\int x^{3}\ln x\,\mathrm{d}x$
  \item $\displaystyle\int q\,\mathrm{e}^{-rq}\,\mathrm{d}q \quad (r>0)$
  \item $\displaystyle\int x\ln(2x)\,\mathrm{d}x$
  \item $\displaystyle\int x^{2}\mathrm{e}^{-x}\,\mathrm{d}x$
        \quad \emph{(tabular)}
  \item $\displaystyle\int (3x-1)\mathrm{e}^{x}\,\mathrm{d}x$
  \item $\displaystyle\int \frac{\ln x}{x^{2}}\,\mathrm{d}x$
  \item $\displaystyle\int x\sqrt{x+1}\,\mathrm{d}x$
        \quad \emph{(parts or substitution)}
  \item $\displaystyle\int x\,2^{x}\,\mathrm{d}x$
  \item $\displaystyle\int x^{2}\mathrm{e}^{2x}\,\mathrm{d}x$
  \item $\displaystyle\int \ln(x+1)\,\mathrm{d}x$
  \item $\displaystyle\int x\,\mathrm{e}^{-2x}\,\mathrm{d}x$
  \item $\displaystyle\int x^{3}\mathrm{e}^{x}\,\mathrm{d}x$
        \quad \emph{(tabular)}
  \item $\displaystyle\int \ln(x^{2})\,\mathrm{d}x$
        \quad \emph{(rewrite first)}
  \item $\displaystyle\int x(\ln x)^{2}\,\mathrm{d}x$
  \item $\displaystyle\int_{0}^{1} x\,\mathrm{e}^{x}\,\mathrm{d}x$
        \quad \emph{(definite)}
  \item $\displaystyle\int_{1}^{\mathrm{e}} \ln x\,\mathrm{d}x$
        \quad \emph{(definite)}
\end{problems}

\section{Solutions}

\begin{solutions}
  \item $u=x$, $v'=\mathrm{e}^{x}$; $u'=1$, $v=\mathrm{e}^{x}$.
        \[
          \int x\mathrm{e}^{x}\,\mathrm{d}x
            = x\mathrm{e}^{x}-\int \mathrm{e}^{x}\,\mathrm{d}x
            = (x-1)\mathrm{e}^{x}+C .
        \]

  \item $u=x$, $v'=\mathrm{e}^{-x}$; $u'=1$, $v=-\mathrm{e}^{-x}$.
        \[
          \int x\mathrm{e}^{-x}\,\mathrm{d}x
            = -x\mathrm{e}^{-x}-\int(-\mathrm{e}^{-x})\,\mathrm{d}x
            = -x\mathrm{e}^{-x}-\mathrm{e}^{-x}+C
            = -(x+1)\mathrm{e}^{-x}+C .
        \]

  \item $u=\ln x$, $v'=x$; $u'=\tfrac1x$, $v=\tfrac{x^{2}}{2}$.
        \[
          \int x\ln x\,\mathrm{d}x
            = \tfrac{x^{2}}{2}\ln x - \int \tfrac{x^{2}}{2}\cdot\tfrac1x\,\mathrm{d}x
            = \tfrac{x^{2}}{2}\ln x - \tfrac{x^{2}}{4}+C .
        \]

  \item $u=\ln x$, $v'=1$; $u'=\tfrac1x$, $v=x$.
        \[
          \int \ln x\,\mathrm{d}x = x\ln x - \int 1\,\mathrm{d}x = x\ln x - x + C .
        \]

  \item $u=x^{2}$, $v'=\mathrm{e}^{x}$; reduce to problem 1.
        \[
          \int x^{2}\mathrm{e}^{x}\,\mathrm{d}x
            = x^{2}\mathrm{e}^{x}-\int 2x\mathrm{e}^{x}\,\mathrm{d}x
            = x^{2}\mathrm{e}^{x}-2(x-1)\mathrm{e}^{x}+C
            = (x^{2}-2x+2)\mathrm{e}^{x}+C .
        \]

  \item $u=\ln x$, $v'=x^{2}$; $v=\tfrac{x^{3}}{3}$.
        \[
          \int x^{2}\ln x\,\mathrm{d}x
            = \tfrac{x^{3}}{3}\ln x - \int \tfrac{x^{3}}{3}\cdot\tfrac1x\,\mathrm{d}x
            = \tfrac{x^{3}}{3}\ln x - \tfrac{x^{3}}{9}+C .
        \]

  \item $u=(\ln x)^{2}$, $v'=1$; $u'=\tfrac{2\ln x}{x}$, $v=x$.
        \[
          \int (\ln x)^{2}\,\mathrm{d}x
            = x(\ln x)^{2}-\int 2\ln x\,\mathrm{d}x
            = x(\ln x)^{2}-2(x\ln x - x)+C
            = x(\ln x)^{2}-2x\ln x+2x+C .
        \]

  \item $u=x$, $v'=\mathrm{e}^{2x}$; $v=\tfrac12\mathrm{e}^{2x}$.
        \[
          \int x\mathrm{e}^{2x}\,\mathrm{d}x
            = \tfrac{x}{2}\mathrm{e}^{2x}-\int \tfrac12\mathrm{e}^{2x}\,\mathrm{d}x
            = \tfrac{x}{2}\mathrm{e}^{2x}-\tfrac14\mathrm{e}^{2x}+C
            = \tfrac14(2x-1)\mathrm{e}^{2x}+C .
        \]

  \item $u=x$, $v'=(x+1)^{8}$; $v=\dfrac{(x+1)^{9}}{9}$.
        \[
          \int x(x+1)^{8}\,\mathrm{d}x
            = \frac{x(x+1)^{9}}{9}-\int \frac{(x+1)^{9}}{9}\,\mathrm{d}x
            = \frac{x(x+1)^{9}}{9}-\frac{(x+1)^{10}}{90}+C .
        \]

  \item $u=\ln x$, $v'=x^{3}$; $v=\tfrac{x^{4}}{4}$.
        \[
          \int x^{3}\ln x\,\mathrm{d}x
            = \tfrac{x^{4}}{4}\ln x - \int \tfrac{x^{4}}{4}\cdot\tfrac1x\,\mathrm{d}x
            = \tfrac{x^{4}}{4}\ln x - \tfrac{x^{4}}{16}+C .
        \]

  \item $u=q$, $v'=\mathrm{e}^{-rq}$; $v=-\tfrac1r\mathrm{e}^{-rq}$.
        \[
          \int q\mathrm{e}^{-rq}\,\mathrm{d}q
            = -\tfrac{q}{r}\mathrm{e}^{-rq}+\tfrac1r\int \mathrm{e}^{-rq}\,\mathrm{d}q
            = -\tfrac{q}{r}\mathrm{e}^{-rq}-\tfrac{1}{r^{2}}\mathrm{e}^{-rq}+C .
        \]
        Equivalently, $-\dfrac{\mathrm{e}^{-rq}(rq+1)}{r^{2}}+C$.

  \item $u=\ln(2x)$, $v'=x$; $u'=\tfrac1x$, $v=\tfrac{x^{2}}{2}$.
        \[
          \int x\ln(2x)\,\mathrm{d}x
            = \tfrac{x^{2}}{2}\ln(2x)-\int \tfrac{x}{2}\,\mathrm{d}x
            = \tfrac{x^{2}}{2}\ln(2x)-\tfrac{x^{2}}{4}+C .
        \]

  \item Tabular. Differentiate $x^{2}$: $x^{2},2x,2,0$. Antidifferentiate
        $\mathrm{e}^{-x}$: $-\mathrm{e}^{-x},\mathrm{e}^{-x},-\mathrm{e}^{-x},\mathrm{e}^{-x}$.
        Sign pattern $+,-,+$:
        \[
          \int x^{2}\mathrm{e}^{-x}\,\mathrm{d}x
            = -x^{2}\mathrm{e}^{-x}-2x\mathrm{e}^{-x}-2\mathrm{e}^{-x}+C
            = -(x^{2}+2x+2)\mathrm{e}^{-x}+C .
        \]
        Check by differentiating:
        $\dfrac{\mathrm{d}}{\mathrm{d}x}\bigl[-(x^{2}+2x+2)\mathrm{e}^{-x}\bigr]
          = -(2x+2)\mathrm{e}^{-x}+(x^{2}+2x+2)\mathrm{e}^{-x}=x^{2}\mathrm{e}^{-x}$.
        \checkmark

  \item Linearity then parts:
        \[
          \int (3x-1)\mathrm{e}^{x}\,\mathrm{d}x
            = 3\int x\mathrm{e}^{x}\,\mathrm{d}x - \int \mathrm{e}^{x}\,\mathrm{d}x
            = 3(x-1)\mathrm{e}^{x}-\mathrm{e}^{x}+C
            = (3x-4)\mathrm{e}^{x}+C .
        \]

  \item $u=\ln x$, $v'=x^{-2}$; $v=-x^{-1}$.
        \[
          \int \frac{\ln x}{x^{2}}\,\mathrm{d}x
            = -\frac{\ln x}{x}+\int \frac{1}{x^{2}}\,\mathrm{d}x
            = -\frac{\ln x}{x}-\frac1x+C
            = -\frac{\ln x+1}{x}+C .
        \]

  \item $u=x$, $v'=(x+1)^{1/2}$; $v=\tfrac23(x+1)^{3/2}$.
        \[
          \int x\sqrt{x+1}\,\mathrm{d}x
            = \tfrac{2x}{3}(x+1)^{3/2}-\int \tfrac23(x+1)^{3/2}\,\mathrm{d}x
            = \tfrac{2x}{3}(x+1)^{3/2}-\tfrac{4}{15}(x+1)^{5/2}+C .
        \]

  \item $u=x$, $v'=2^{x}$; $v=\dfrac{2^{x}}{\ln 2}$.
        \[
          \int x\,2^{x}\,\mathrm{d}x
            = \frac{x\,2^{x}}{\ln 2}-\int \frac{2^{x}}{\ln 2}\,\mathrm{d}x
            = \frac{x\,2^{x}}{\ln 2}-\frac{2^{x}}{(\ln 2)^{2}}+C .
        \]

  \item Tabular. Differentiate $x^{2}$: $x^{2},2x,2,0$. Antidifferentiate
        $\mathrm{e}^{2x}$: $\tfrac12\mathrm{e}^{2x},\tfrac14\mathrm{e}^{2x},
        \tfrac18\mathrm{e}^{2x},\tfrac{1}{16}\mathrm{e}^{2x}$. Signs $+,-,+$:
        \[
          \int x^{2}\mathrm{e}^{2x}\,\mathrm{d}x
            = \tfrac{x^{2}}{2}\mathrm{e}^{2x}-\tfrac{x}{2}\mathrm{e}^{2x}
              +\tfrac14\mathrm{e}^{2x}+C
            = \tfrac{\mathrm{e}^{2x}}{4}(2x^{2}-2x+1)+C .
        \]

  \item $u=\ln(x+1)$, $v'=1$; $u'=\tfrac{1}{x+1}$, $v=x$. To avoid an
        awkward residual, write $v=x+1$ instead (legal, since $v$ is any
        antiderivative of $1$):
        \[
          \int \ln(x+1)\,\mathrm{d}x
            = (x+1)\ln(x+1)-\int 1\,\mathrm{d}x
            = (x+1)\ln(x+1)-x+C .
        \]

  \item $u=x$, $v'=\mathrm{e}^{-2x}$; $v=-\tfrac12\mathrm{e}^{-2x}$.
        \[
          \int x\mathrm{e}^{-2x}\,\mathrm{d}x
            = -\tfrac{x}{2}\mathrm{e}^{-2x}+\tfrac12\int \mathrm{e}^{-2x}\,\mathrm{d}x
            = -\tfrac{x}{2}\mathrm{e}^{-2x}-\tfrac14\mathrm{e}^{-2x}+C
            = -\tfrac{\mathrm{e}^{-2x}}{4}(2x+1)+C .
        \]

  \item Tabular. Differentiate $x^{3}$: $x^{3},3x^{2},6x,6,0$.
        Antidifferentiate $\mathrm{e}^{x}$: $\mathrm{e}^{x}$ throughout.
        Signs $+,-,+,-$:
        \[
          \int x^{3}\mathrm{e}^{x}\,\mathrm{d}x
            = (x^{3}-3x^{2}+6x-6)\mathrm{e}^{x}+C .
        \]

  \item Rewrite: $\ln(x^{2})=2\ln|x|$, so
        \[
          \int 2\ln|x|\,\mathrm{d}x
            = 2(x\ln|x|-x)+C = 2x\ln|x|-2x+C .
        \]

  \item $u=(\ln x)^{2}$, $v'=x$; $u'=\dfrac{2\ln x}{x}$, $v=\dfrac{x^{2}}{2}$.
        \[
          \int x(\ln x)^{2}\,\mathrm{d}x
            = \tfrac{x^{2}}{2}(\ln x)^{2}-\int x\ln x\,\mathrm{d}x
            = \tfrac{x^{2}}{2}(\ln x)^{2}-\tfrac{x^{2}}{2}\ln x+\tfrac{x^{2}}{4}+C ,
        \]
        using problem 3 for the residual.

  \item By problem 1 the antiderivative is $(x-1)\mathrm{e}^{x}$. Evaluate:
        \[
          \int_{0}^{1} x\mathrm{e}^{x}\,\mathrm{d}x
            = \bigl[(x-1)\mathrm{e}^{x}\bigr]_{0}^{1}
            = 0-(-1)=1 .
        \]

  \item Antiderivative $x\ln x - x$.
        \[
          \int_{1}^{\mathrm{e}} \ln x\,\mathrm{d}x
            = \bigl[x\ln x - x\bigr]_{1}^{\mathrm{e}}
            = (\mathrm{e}\cdot 1-\mathrm{e})-(1\cdot 0-1)
            = 0-(-1)=1 .
        \]
\end{solutions}

\section{Additional Problems}

Compute each integral using parts (possibly more than once).

\begin{problems}
  \item $\displaystyle\int x\,\mathrm{e}^{3x}\,\mathrm{d}x$
  \item $\displaystyle\int x\,\mathrm{e}^{-3x}\,\mathrm{d}x$
  \item $\displaystyle\int x^{2}\mathrm{e}^{-2x}\,\mathrm{d}x$
        \quad \emph{(tabular)}
  \item $\displaystyle\int x^{4}\ln x\,\mathrm{d}x$
  \item $\displaystyle\int x\ln(3x)\,\mathrm{d}x$
  \item $\displaystyle\int (2x+5)\mathrm{e}^{x}\,\mathrm{d}x$
  \item $\displaystyle\int \ln(2x+1)\,\mathrm{d}x$
  \item $\displaystyle\int \frac{\ln x}{x^{3}}\,\mathrm{d}x$
  \item $\displaystyle\int x\,3^{x}\,\mathrm{d}x$
  \item $\displaystyle\int x^{3}\mathrm{e}^{-x}\,\mathrm{d}x$
        \quad \emph{(tabular)}
  \item $\displaystyle\int q^{2}\mathrm{e}^{-rq}\,\mathrm{d}q \quad (r>0)$
  \item $\displaystyle\int \sqrt{x}\ln x\,\mathrm{d}x$
  \item $\displaystyle\int x\sqrt{2x+1}\,\mathrm{d}x$
        \quad \emph{(parts or substitution)}
  \item $\displaystyle\int_{0}^{1} x\,\mathrm{e}^{2x}\,\mathrm{d}x$
        \quad \emph{(definite)}
  \item $\displaystyle\int_{1}^{\mathrm{e}} x^{2}\ln x\,\mathrm{d}x$
        \quad \emph{(definite)}
\end{problems}

\section{Additional Solutions}

\begin{solutions}
  \item $u=x$, $v'=\mathrm{e}^{3x}$; $u'=1$, $v=\tfrac13\mathrm{e}^{3x}$.
        \[
          \int x\mathrm{e}^{3x}\,\mathrm{d}x
            = \tfrac{x}{3}\mathrm{e}^{3x}-\int \tfrac13\mathrm{e}^{3x}\,\mathrm{d}x
            = \tfrac{x}{3}\mathrm{e}^{3x}-\tfrac19\mathrm{e}^{3x}+C
            = \tfrac{\mathrm{e}^{3x}}{9}(3x-1)+C .
        \]

  \item $u=x$, $v'=\mathrm{e}^{-3x}$; $v=-\tfrac13\mathrm{e}^{-3x}$.
        \[
          \int x\mathrm{e}^{-3x}\,\mathrm{d}x
            = -\tfrac{x}{3}\mathrm{e}^{-3x}+\tfrac13\int \mathrm{e}^{-3x}\,\mathrm{d}x
            = -\tfrac{x}{3}\mathrm{e}^{-3x}-\tfrac19\mathrm{e}^{-3x}+C
            = -\tfrac{\mathrm{e}^{-3x}}{9}(3x+1)+C .
        \]

  \item Tabular. Differentiate $x^{2}$: $x^{2},2x,2,0$. Antidifferentiate
        $\mathrm{e}^{-2x}$: $-\tfrac12\mathrm{e}^{-2x},\tfrac14\mathrm{e}^{-2x},
        -\tfrac18\mathrm{e}^{-2x}$. Signs $+,-,+$:
        \[
          \int x^{2}\mathrm{e}^{-2x}\,\mathrm{d}x
            = -\tfrac{x^{2}}{2}\mathrm{e}^{-2x}-\tfrac{x}{2}\mathrm{e}^{-2x}
              -\tfrac14\mathrm{e}^{-2x}+C
            = -\tfrac{\mathrm{e}^{-2x}}{4}(2x^{2}+2x+1)+C .
        \]

  \item $u=\ln x$, $v'=x^{4}$; $v=\tfrac{x^{5}}{5}$.
        \[
          \int x^{4}\ln x\,\mathrm{d}x
            = \tfrac{x^{5}}{5}\ln x - \int \tfrac{x^{5}}{5}\cdot\tfrac1x\,\mathrm{d}x
            = \tfrac{x^{5}}{5}\ln x - \tfrac{x^{5}}{25}+C .
        \]

  \item $u=\ln(3x)$, $v'=x$; $u'=\tfrac1x$, $v=\tfrac{x^{2}}{2}$.
        \[
          \int x\ln(3x)\,\mathrm{d}x
            = \tfrac{x^{2}}{2}\ln(3x)-\int \tfrac{x}{2}\,\mathrm{d}x
            = \tfrac{x^{2}}{2}\ln(3x)-\tfrac{x^{2}}{4}+C .
        \]

  \item Linearity then parts:
        \[
          \int (2x+5)\mathrm{e}^{x}\,\mathrm{d}x
            = 2\int x\mathrm{e}^{x}\,\mathrm{d}x + 5\int \mathrm{e}^{x}\,\mathrm{d}x
            = 2(x-1)\mathrm{e}^{x}+5\mathrm{e}^{x}+C
            = (2x+3)\mathrm{e}^{x}+C .
        \]

  \item $u=\ln(2x+1)$, $v'=1$; $u'=\tfrac{2}{2x+1}$. Take $v=x+\tfrac12$
        (any antiderivative of $1$). Then
        \[
          \int \ln(2x+1)\,\mathrm{d}x
            = \tfrac{2x+1}{2}\ln(2x+1)-\int \tfrac{(2x+1)\cdot 2}{2(2x+1)}\,\mathrm{d}x
            = \tfrac{2x+1}{2}\ln(2x+1)-x+C .
        \]

  \item $u=\ln x$, $v'=x^{-3}$; $u'=\tfrac1x$, $v=-\tfrac12 x^{-2}$.
        \[
          \int \frac{\ln x}{x^{3}}\,\mathrm{d}x
            = -\frac{\ln x}{2x^{2}}+\int \frac{1}{2x^{3}}\,\mathrm{d}x
            = -\frac{\ln x}{2x^{2}}-\frac{1}{4x^{2}}+C
            = -\frac{2\ln x+1}{4x^{2}}+C .
        \]

  \item $u=x$, $v'=3^{x}$; $v=\dfrac{3^{x}}{\ln 3}$.
        \[
          \int x\,3^{x}\,\mathrm{d}x
            = \frac{x\,3^{x}}{\ln 3}-\int \frac{3^{x}}{\ln 3}\,\mathrm{d}x
            = \frac{x\,3^{x}}{\ln 3}-\frac{3^{x}}{(\ln 3)^{2}}+C .
        \]

  \item Tabular. Differentiate $x^{3}$: $x^{3},3x^{2},6x,6,0$.
        Antidifferentiate $\mathrm{e}^{-x}$:
        $-\mathrm{e}^{-x},\mathrm{e}^{-x},-\mathrm{e}^{-x},\mathrm{e}^{-x}$.
        Signs $+,-,+,-$:
        \[
          \int x^{3}\mathrm{e}^{-x}\,\mathrm{d}x
            = -x^{3}\mathrm{e}^{-x}-3x^{2}\mathrm{e}^{-x}-6x\mathrm{e}^{-x}
              -6\mathrm{e}^{-x}+C
            = -(x^{3}+3x^{2}+6x+6)\mathrm{e}^{-x}+C .
        \]

  \item $u=q^{2}$, $v'=\mathrm{e}^{-rq}$; $v=-\tfrac1r\mathrm{e}^{-rq}$.
        \[
          \int q^{2}\mathrm{e}^{-rq}\,\mathrm{d}q
            = -\tfrac{q^{2}}{r}\mathrm{e}^{-rq}+\tfrac2r\int q\mathrm{e}^{-rq}\,\mathrm{d}q .
        \]
        Using problem 11 of the practice problems,
        $\int q\mathrm{e}^{-rq}\,\mathrm{d}q = -\tfrac{\mathrm{e}^{-rq}(rq+1)}{r^{2}}$, so
        \[
          \int q^{2}\mathrm{e}^{-rq}\,\mathrm{d}q
            = -\tfrac{q^{2}}{r}\mathrm{e}^{-rq}-\tfrac{2\mathrm{e}^{-rq}(rq+1)}{r^{3}}+C
            = -\tfrac{\mathrm{e}^{-rq}}{r^{3}}\bigl(r^{2}q^{2}+2rq+2\bigr)+C .
        \]

  \item $u=\ln x$, $v'=x^{1/2}$; $u'=\tfrac1x$, $v=\tfrac23 x^{3/2}$.
        \[
          \int \sqrt{x}\ln x\,\mathrm{d}x
            = \tfrac23 x^{3/2}\ln x - \int \tfrac23 x^{1/2}\,\mathrm{d}x
            = \tfrac23 x^{3/2}\ln x - \tfrac49 x^{3/2}+C .
        \]

  \item $u=x$, $v'=(2x+1)^{1/2}$; $v=\tfrac13(2x+1)^{3/2}$.
        \[
          \int x\sqrt{2x+1}\,\mathrm{d}x
            = \tfrac{x}{3}(2x+1)^{3/2}-\tfrac13\int (2x+1)^{3/2}\,\mathrm{d}x
            = \tfrac{x}{3}(2x+1)^{3/2}-\tfrac{1}{15}(2x+1)^{5/2}+C .
        \]

  \item Antiderivative of $x\mathrm{e}^{2x}$ is $\tfrac14(2x-1)\mathrm{e}^{2x}$
        (problem 8 of the practice problems). Evaluate:
        \[
          \int_{0}^{1} x\mathrm{e}^{2x}\,\mathrm{d}x
            = \bigl[\tfrac14(2x-1)\mathrm{e}^{2x}\bigr]_{0}^{1}
            = \tfrac{\mathrm{e}^{2}}{4}-\bigl(-\tfrac14\bigr)
            = \tfrac{\mathrm{e}^{2}+1}{4} .
        \]

  \item Antiderivative of $x^{2}\ln x$ is $\tfrac{x^{3}}{3}\ln x-\tfrac{x^{3}}{9}$
        (problem 6 of the practice problems). Evaluate:
        \[
          \int_{1}^{\mathrm{e}} x^{2}\ln x\,\mathrm{d}x
            = \bigl[\tfrac{x^{3}}{3}\ln x-\tfrac{x^{3}}{9}\bigr]_{1}^{\mathrm{e}}
            = \bigl(\tfrac{\mathrm{e}^{3}}{3}-\tfrac{\mathrm{e}^{3}}{9}\bigr)
              -\bigl(0-\tfrac19\bigr)
            = \tfrac{2\mathrm{e}^{3}}{9}+\tfrac19
            = \tfrac{2\mathrm{e}^{3}+1}{9} .
        \]
\end{solutions}
